\section{Related Work}
\label{sec:relatedwork}

This section reviews prior work along three axes that are relevant
to the proposed prototype: smart product identification methods,
sensor-based smart-refrigerator systems, and visual food-quality
analysis. We then position the proposed module against these axes.

\subsection{Smart Product Identification Methods}

Product identification is a core component of any inventory-monitoring
system. The literature has converged on three families of approaches:
RFID, barcode-based, and camera-assisted methods.

\textbf{RFID-based systems} eliminate manual data entry and can
identify items even when partially occluded, which is attractive for
refrigerator-style enclosures~\cite{abhiram2025iotmonitor}. The
practical limitations are well documented: every monitored product
must carry an RFID tag, a reader has to be installed inside the
enclosure, and the per-unit cost of tags and reader hardware is
non-trivial for household-scale prototypes. In Türkiye's retail
context, refrigerated household goods are rarely pre-tagged with RFID,
which makes the approach impractical without manual re-tagging by the
end user --- arguably worse than the manual-entry problem it was
meant to solve.

\textbf{Barcode-based systems} are cheap and ubiquitous because
barcodes are already printed on most packaged products. However,
barcodes require careful alignment with a scanner, are sensitive to
crumpling and condensation, and frequently sit on the package face
that ends up against the back of the refrigerator
shelf~\cite{pelobello2026chickout}. As a result, daily-use barcode
systems either tend to be abandoned by their users or require an
additional manual "rescan" step that defeats the purpose of
automation.

\textbf{Camera-based recognition} methods identify products through
on-shelf images using either generic image classifiers or trained
domain-specific models~\cite{ericyang2022smartfridge,vadlamudi2025fridgesense}.
These approaches can in principle reach high accuracy but pay for it
in three ways: they require large training datasets that often do not
exist for regional product brands, they consume significant compute
either on the edge device or via a cloud round-trip, and their
behavior degrades quickly when lighting or product packaging changes.
Several authors note that classifier confidence drops substantially
when fewer than ten frames per product are available for training,
which is the realistic case for household
deployments~\cite{vadlamudi2025fridgesense}.

A summary comparison of these methods alongside the proposed
QR-based approach is given in Table~\ref{tab:identification-methods}.

\begin{table*}[t]
\caption{Comparison of product identification methods in smart inventory systems.}
\label{tab:identification-methods}
\centering
\renewcommand{\arraystretch}{1.2}
\begin{tabular}{@{}lcccc@{}}
\toprule
\textbf{Method} & \textbf{Automation} &
\textbf{Required hardware} & \textbf{Per-unit cost} &
\textbf{Data capacity per identifier} \\
\midrule
RFID                  & High        & RFID tags + reader & High & Medium \\
Barcode               & Low         & Barcode scanner    & Low  & Low \\
QR-based identification & Medium--High & Camera module    & Low  & High \\
\textbf{Proposed (QR + ESP32-CAM)} & \textbf{High} &
\textbf{ESP32-CAM} & \textbf{Low} & \textbf{High} \\
\bottomrule
\end{tabular}
\end{table*}

\subsection{Sensor-Based Smart-Refrigerator Systems}

A second family of related work concerns sensor instrumentation of the
storage enclosure itself. Goel \emph{et al.} present a prototype
combining a small camera with weight sensors to identify shopping
needs and partially log inventory~\cite{goel2023iotfridge}. Sarkar
\emph{et al.}~\cite{sarkar2024proactive} extend this idea with
proactive notifications based on consumption patterns. Shetty and
Shetty propose an IoT-monitored food-storage system that emphasizes
environmental conditions~\cite{shetty2024monitoring}. Nandini
\emph{et al.}~\cite{nandini2023userconvenience} explore the same
problem domain through the lens of usability improvements to the
refrigerator interface.

Despite this body of work, most published prototypes optimize for one
sensing modality at a time. Combinations of weight tracking,
environmental sensing, and automatic product identification within a
single low-cost reference design are rare in the open
literature~\cite{sarawathi2025spoilage}. The present work is positioned
explicitly to fill this gap: a single Arduino~Uno node coordinating
weight, temperature, humidity, and gas-trend monitoring; an auxiliary
ESP32-CAM node for QR scanning and visual ripeness analysis; and a
unified mobile interface that exposes every measurement and event in
real time.

\subsection{Visual Food-Quality Analysis}

A third strand of related work addresses the analysis of visual
spoilage indicators, of which banana browning is a common benchmark
because the surface color of bananas changes predictably with
ripeness. The dominant approach in the recent literature is
convolutional neural network (CNN) classification using transfer
learning from large-scale visual
backbones~\cite{vadlamudi2025fridgesense}. CNNs can reach high
accuracy on representative test sets but require significant
training data, edge-side compute, or a cloud connection --- none of
which are attractive for low-budget household prototypes that must
also operate on intermittent Wi-Fi.

The proposed module instead uses a deliberately simple, fully
interpretable pixel-level RGB analysis to estimate the visible
"spotting percentage" of a banana surface. This approach trades a
small accuracy gap (relative to CNN-based detectors) for transparency,
zero training data requirements, and execution times of a few
milliseconds on a midrange smartphone. It is intended as an
\emph{honest baseline} that can be replaced by a CNN model in future
work, while still providing useful information out of the box.

\subsection{Positioning of the Proposed Module}

In summary, the proposed Smart Inventory and Automated Food Tracking
Module differs from prior art in four respects:

\begin{enumerate}
  \item It \emph{combines} four sensing modalities --- weight,
        temperature/humidity, gas trend, and visual --- inside a single
        low-cost cabinet, rather than focusing on a single dimension.
  \item It uses \emph{QR-based} product registration with an ESP32-CAM,
        which is markedly cheaper than RFID and more reliable than
        barcode scanning in cabinet conditions.
  \item It exposes \emph{real-time mobile} visibility through a
        cross-platform Flutter application backed by a small Python
        bridge, removing the need for a Wi-Fi-capable
        microcontroller on the sensor side.
  \item It provides a \emph{transparent, interpretable} visual-ripeness
        baseline that does not require machine-learning infrastructure,
        which we believe is the right starting point for a reference
        design intended for educational and reproducibility use.
\end{enumerate}

The remaining sections describe the system design, implementation, and
qualitative evaluation in detail.
